\tikzset{node distance=0pt,
  promtitle/.style  = {draw=tugb!70, fill=tugb!15, rounded corners=5pt, inner xsep=10pt, inner ysep=6pt, font=\small\bfseries\sffamily, text=tugb!90!black, align=left},
  promsechdr/.style = {draw=none, fill=tugb!12, rounded corners=0pt, inner xsep=5pt, inner ysep=3pt, font=\scriptsize\bfseries\sffamily, text=tugb!80!black, anchor=west},
  promrole/.style   = {draw=none, fill=tugb!10, inner sep=1mm, rounded corners=0pt, align=left, font=\scriptsize\ttfamily},
  promapi/.style    = {draw=none, fill=darkgreencode!3, inner sep=1mm, rounded corners=0pt, align=left, font=\scriptsize\ttfamily},
  promphi/.style    = {draw=none, fill=darkgreencode!9, rounded corners=0pt, inner sep=1mm, align=left, font=\scriptsize\ttfamily},
  prompta/.style    = {draw=none, fill=darkgreencode!13, rounded corners=0pt, inner sep=1mm, align=left, font=\scriptsize\ttfamily},
  prompfi/.style    = {draw=none, fill=darkgreencode!17, rounded corners=0pt, inner sep=1mm, align=left, font=\scriptsize\ttfamily},
  prompexample/.style    = {draw=none, fill=darkgreencode!21, rounded corners=0pt, inner sep=1mm, align=left, font=\scriptsize\ttfamily},
  promex/.style     = {draw=gray!35, fill=darkgreencode!5, rounded corners=0pt, inner sep=8pt, align=left, font=\scriptsize\ttfamily},
  promtask/.style   = {draw=none, fill=tug!6, rounded corners=0pt, inner sep=1mm, align=left, font=\scriptsize\ttfamily},
}
\begin{tikzpicture}

%%─────────────────────────────────────────────────────────────────────────────
%% [1] ROLE
%%─────────────────────────────────────────────────────────────────────────────
\node[promrole, text width=0.5\textwidth] (role) { \textcolor{tug}{\textbf{{[1]\;ROLE:}}} You are an expert C++17 developer for the \textbf{SOLEX} logical query rewriter. Given a formal rule specification $r^* = (\phi_{r^*},\,\tau_{r^*},\,\mathcal{K}_{r^*})$, implement the corresponding \emph{Rule} subclass with \emph{Matches()} and \emph{Apply()}.};

%%─────────────────────────────────────────────────────────────────────────────
%% [2] RULE API
%%─────────────────────────────────────────────────────────────────────────────
\node[promapi, text width=0.5\textwidth, below=of role] (code) {%
 \textcolor{tug}{\textbf{{[2]\;RULE API:}}} \\
namespace solex::rules \{\\
\phantom{x}class \textcolor{jb-class}{\bfseries Rule} \{\\
\phantom{xx}\textcolor{jb-comment}{// \textbf{Matches}(): returns true iff $\phi_{r^*}$ holds on op; called on every node}\\
\phantom{xx}\textcolor{jb-keyword}{virtual} bool \textcolor{jb-function}{\bfseries Matches}(\textcolor{jb-keyword}{const} \textcolor{jb-class}{ir::LogicalOperator}\& op) \textcolor{jb-keyword}{const} = \textcolor{jb-number}{0};\\[3pt]
\phantom{xx}\textcolor{jb-comment}{// \textbf{Apply}(): transforms plan per $\tau_{r^*}$; return std::move(op) unchanged if inapplicable}\\
\phantom{xx}\textcolor{jb-keyword}{virtual} \textcolor{jb-class}{ir::LogicalOperatorPtr} \textcolor{jb-function}{\bfseries Apply}(\textcolor{jb-class}{ir::LogicalOperatorPtr} op) \textcolor{jb-keyword}{const} = \textcolor{jb-number}{0};\\
\phantom{x}\};\phantom{x}\}\\[4pt]
\textcolor{jb-comment}{// Key types:}\;\textcolor{jb-class}{LogicalOperatorType}\texttt{::}\{\textcolor{jb-field}{AGGREGATE}\;|\;\textcolor{jb-field}{JOIN}\;|\;\textcolor{jb-field}{SCAN}\;|\;\textcolor{jb-field}{FILTER}\;|\;\textcolor{jb-field}{WITH}\;|\;\textcolor{jb-field}{PROJECT}\;\ldots\}\\
\textcolor{jb-comment}{// Key types:}\;\textcolor{jb-class}{JoinType}\texttt{::}\{\textcolor{jb-field}{INNER}\;|\;\textcolor{jb-field}{LEFT}\;|\;\textcolor{jb-field}{RIGHT}\;|\;\textcolor{jb-field}{FULL}\;|\;\textcolor{jb-field}{CROSS}\;\ldots\}\\
\textcolor{jb-comment}{// Accessors:}\;\texttt{op.type()}\;\;\texttt{op.Cast<T>()}\;\;\texttt{op.children()}\;\;\texttt{op.Clone()}\\
\textcolor{jb-comment}{// AggregateOperator:}\;\texttt{.group\_by()}\;\;\texttt{.aggregates()}\\
\textcolor{jb-comment}{// JoinOperator:}\;\texttt{.join\_type()}\;\;\texttt{.predicates()}\\
\textcolor{jb-comment}{// ScanOperator:}\;\texttt{.table\_name()}\;\;\texttt{.alias()}};

%%─────────────────────────────────────────────────────────────────────────────
%% [3] RULE SPECIFICATION  r* = (phi, tau, K)
%%─────────────────────────────────────────────────────────────────────────────
%% phi_r* box (left)
\node[promphi, text width=0.5\textwidth, below=of code, xshift=0pt] (phi_box){%
  \textcolor{tug}{\textbf{{[3]\;$\phi_{r^*}$\enspace Applicability Condition:}}}\\[3pt]
  $\bullet$\enspace\textbf{Root}:\;\texttt{AGG[COUNT(*)]} with empty \texttt{GROUP BY}\enspace(no grouping columns)\\[2pt]
  $\bullet$\enspace\textbf{Subtree}:\;$\geq 3$ scan nodes connected by \texttt{INNER JOIN}s forming a star\\[2pt]
  $\bullet$\enspace\textbf{FK structure}:\;all inter-table join predicates are equalities $T_i.fk_i = T_{anc}.pk$ for a single anchor table $T_{anc}$\\[2pt]
  $\bullet$\enspace Every non-anchor $T_i$ links directly to $T_{anc}$\enspace(star topology; no transitive or hub-bypass joins)
  };

%% tau_r* + K_r* box (right)
\node[prompta, text width=0.5\textwidth, below=of phi_box] (tau_box) {%
  \textcolor{tug}{\textbf{[4] $\tau_{r^*}$\enspace Rewrite Transformation:}}\\[3pt]
  $\bullet$\enspace For each non-anchor branch $T_i$: build $\textit{CTE}_i$\;=\;\texttt{AGG[COUNT(*)} \texttt{GROUP BY}~$fk_i$\texttt{]}  over \texttt{SCAN[}$T_i$\texttt{]} + per-table filters\\[2pt]
  $\bullet$\enspace Join $T_{anc}$ to every $\textit{CTE}_i$ via $T_{anc}.pk = \textit{CTE}_i.fk_i$ (\texttt{INNER})\\[2pt]
  $\bullet$\enspace Replace root with: \texttt{AGG[SUM(}$b_1.n\!\times\!\cdots\!\times\!b_N.n$\texttt{)]}};  

\node[prompfi, text width=0.5\textwidth, below=of tau_box] (FI) {%
  \textcolor{tug}{\textbf{[5] $\mathcal{K}_{r^*}$\enspace Fragment Keys:}} 
  \textcolor{jb-keyword}{$\blacksquare$}\;$\mathcal{K}_1$\enspace\texttt{AGG(FILTER)}
  \textcolor{jb-keyword}{$\blacksquare$}\;$\mathcal{K}_2$\enspace\texttt{AGG(FILTER(JOIN[INNER]))}
  \textcolor{jb-keyword}{$\blacksquare$}\;$\mathcal{K}_3$\enspace\texttt{AGG(JOIN[INNER])}
  \textcolor{jb-keyword}{$\blacksquare$}\;$\mathcal{K}_4$\enspace\texttt{AGG(FILTER(JOIN[INNER](JOIN,SCAN)))}
  \textcolor{jb-keyword}{$\blacksquare$}\;$\mathcal{K}_5$\enspace\texttt{AGG(JOIN[INNER](JOIN,SCAN))}
  };  

\node[prompexample, text width=0.5\textwidth, below=of FI] (example) {%
  \textcolor{tug}{\textbf{[6] Example: Distillation Seed from Which $r^*$ Was Extracted:}}\\[3pt] 
  <Original SQL Query ($Q$)>, <Rewritten SQL Query ($Q''$)>};   
%%─────────────────────────────────────────────────────────────────────────────
%% [5] TASK
%%─────────────────────────────────────────────────────────────────────────────
\node[promtask, text width=0.5\textwidth, below=of example]{%
  \textcolor{tug}{\textbf{[7] Task:}} Implement Matches() and Apply() for the rule above. Matches() must return True if and only if the operator satisfies $\phi_{r^*}$ exactly as specified. Apply() must produce the rewritten plan according to $\tau_{r^*}$ and return a valid LogicalOperatorPtr. \textbf{Output only compilable C++17 code} for the two method bodies.\enspace No explanation.\enspace No markdown fences.};

\end{tikzpicture}